{{دیوان
| عنوان = شبلیؔ و حالیؔ
| ماخذ = ویکی ماخذ
| ماخذ_لنک = https://ur.wikisource.org/wiki/%D8%B4%D8%A8%D9%84%DB%8C_%D9%88%D8%AD%D8%A7%D9%84%DB%8C
| مقابلہ = اقبال سائبر لائبریری (اقبال اکادمی پاکستان)
| مقابلہ_لنک = https://www.iqbalcyberlibrary.net
}}
<poem>
مسلم سے ایک روز یہ اقبالؔ نے کہا
دیوانِ جُزو و کُل میں ہے تیرا وجود فرد

تیرے سرودِ رفتہ کے نغمے علومِ نَو
تہذیب تیرے قافلہ ہائے کُہن کی گرد

پتھّر ہے اس کے واسطے موجِ نسیم بھی
نازک بہت ہے آئنۂ آبرُوئے مرد

مردانِ کار، ڈھُونڈ کے اسبابِ حادثات
کرتے ہیں چارۂ ستَمِ چرخِ لاجورد

پُوچھ اُن سے جو چمن کے ہیں دیرینہ رازدار
کیونکر ہوئی خزاں ترے گُلشن سے ہم نبرد

مسلم مرے کلام سے بے تاب ہوگیا
غمّاز ہوگئی غمِ پنہاں کی آہِ سرد

کہنے لگا کہ دیکھ تو کیفیّتِ خزاں
اوراق ہو گئے شجَرِ زندگی کے زرد

خاموش ہو گئے چمَنِستاں کے رازدار
سرمایۂ گداز تھی جن کی نوائے درد

شبلیؔ کو رو رہے تھے ابھی اہلِ گلستاں
حالیؔ بھی ہوگیا سوئے فردوس رہ نورد

“اکنوں کرا دماغ کہ پُرسد زباغباں
بُلبل چہ گفت و گُل چہ شنید و صبا چہ کرد”
</poem>
